% GENERATED FILE. Edit canonical lessons, then run npm run generate:books.
% canonical-source-hash: fnv1a64:52eb77700b8959a6
% canonical-lessons: BN-C42-hat, BN-C42-chul, BN-C42-dat, BN-C42-gola, BN-C42-pet

\chapter{Hand, Hair, Tooth, Throat, Stomach}
\label{ch:bn-hand-hair-tooth-throat-stomach}
\hlchaptermodality{42}

\begin{chapteropening}
\textbf{By the end of this chapter:} \emph{I can say a hand, hair, a tooth, the throat and the stomach.}

Complete the last lesson of chapter 42: I can say a hand, hair, a tooth, the throat and the stomach.
\end{chapteropening}

\section[\texorpdfstring{\bn{হাত} (hāt)}{হাত (hāt)}]{\bn{হাত} (hāt) — a hand}

\label{lesson:BN-C42-hat}



\begin{quote}
Before the new one: say the Bengali for first, then the Bengali for fifth.
\end{quote}

\subsection*{You'll want to know: \bn{হাত}}
\textbf{\bn{হাত}} — \emph{hāt} — \textquotedblleft{}a hand\textquotedblright{}.

\begin{cousinweb}
\textbf{\bn{হাত}} (\emph{hāt}) is a tadbhava, a word inherited from Sanskrit and reshaped by sound change. It comes from \emph{hasta}, \textquotedblleft{}hand\textquotedblright{} (CDIAL 14024).
\end{cousinweb}

\begin{grammarlens}[title={What you've built}]
One more word for this chapter.
\end{grammarlens}

\subsection*{Your turn}
\begin{itemize}
\raggedright
  \item \emph{Say it:} \emph{hāt}
  \item \emph{Say it:} \emph{hāt}, once more
  \item \emph{Say it:} say \emph{hāt}
  \item \emph{Recall:} say the Bengali for first, then the Bengali for fifth, then say \emph{hāt} again
\end{itemize}

\begin{tcolorbox}[breakable,colback=teal!4,colframe=teal!35!black,title={Before you move on}]
What does \bn{হাত} mean? (\textquotedblleft{}A hand\textquotedblright{}.) Say it once more.
\end{tcolorbox}
\section[\texorpdfstring{\bn{চুল} (chul)}{চুল (chul)}]{\bn{চুল} (chul) — hair}

\label{lesson:BN-C42-chul}



\begin{quote}
Before the new one: say the Bengali for fifth, then the Bengali for a hand.
\end{quote}

\subsection*{You'll want to know: \bn{চুল}}
\textbf{\bn{চুল}} — \emph{chul} — \textquotedblleft{}hair\textquotedblright{}.

\begin{cousinweb}
\textbf{\bn{চুল}} (\emph{chul}) is a tadbhava, a word inherited from Sanskrit and reshaped by sound change. It comes from \emph{cūḍa}, \textquotedblleft{}topknot on the head\textquotedblright{} (CDIAL 4883). CDIAL takes \emph{cūḍa} from Dravidian, while noting that it may instead belong with a group of words for \textquotedblleft{}hair\textquotedblright{} traced to Munda.
\end{cousinweb}

\begin{grammarlens}[title={What you've built}]
One more word for this chapter.
\end{grammarlens}

\subsection*{Your turn}
\begin{itemize}
\raggedright
  \item \emph{Say it:} \emph{chul}
  \item \emph{Say it:} \emph{chul}, once more
  \item \emph{Say it:} say \emph{chul}
  \item \emph{Recall:} say the Bengali for fifth, then the Bengali for a hand, then say \emph{chul} again
\end{itemize}

\begin{tcolorbox}[breakable,colback=teal!4,colframe=teal!35!black,title={Before you move on}]
What does \bn{চুল} mean? (\textquotedblleft{}Hair\textquotedblright{}.) Say it once more.
\end{tcolorbox}
\section[\texorpdfstring{\bn{দাঁত} (dā̃t)}{দাঁত (dā̃t)}]{\bn{দাঁত} (dā̃t) — a tooth}

\label{lesson:BN-C42-dat}



\begin{quote}
Before the new one: say the Bengali for a hand, then the Bengali for hair.
\end{quote}

\subsection*{You'll want to know: \bn{দাঁত}}
\textbf{\bn{দাঁত}} — \emph{dā̃t} — \textquotedblleft{}a tooth\textquotedblright{}.

\begin{cousinweb}
\textbf{\bn{দাঁত}} (\emph{dāṁt}) is a tadbhava, a word inherited from Sanskrit and reshaped by sound change. It comes from \emph{danta}, \textquotedblleft{}tooth\textquotedblright{} (CDIAL 6152).
\end{cousinweb}

\begin{grammarlens}[title={What you've built}]
One more word for this chapter.
\end{grammarlens}

\subsection*{Your turn}
\begin{itemize}
\raggedright
  \item \emph{Say it:} \emph{dā̃t}
  \item \emph{Say it:} \emph{dā̃t}, once more
  \item \emph{Say it:} say \emph{dā̃t}
  \item \emph{Recall:} say the Bengali for a hand, then the Bengali for hair, then say \emph{dā̃t} again
\end{itemize}

\begin{tcolorbox}[breakable,colback=teal!4,colframe=teal!35!black,title={Before you move on}]
What does \bn{দাঁত} mean? (\textquotedblleft{}A tooth\textquotedblright{}.) Say it once more.
\end{tcolorbox}
\section[\texorpdfstring{\bn{গলা} (gôlā)}{গলা (gôlā)}]{\bn{গলা} (gôlā) — the throat}

\label{lesson:BN-C42-gola}



\begin{quote}
Before the new one: say the Bengali for hair, then the Bengali for a tooth.
\end{quote}

\subsection*{You'll want to know: \bn{গলা}}
\textbf{\bn{গলা}} — \emph{gôlā} — \textquotedblleft{}the throat\textquotedblright{}.

\begin{cousinweb}
\textbf{\bn{গলা}} (\emph{gôlā}) is a tadbhava, a word inherited from Sanskrit and reshaped by sound change. It comes from \emph{gala}, \textquotedblleft{}throat, neck\textquotedblright{} (CDIAL 4070).
\end{cousinweb}

\begin{grammarlens}[title={What you've built}]
One more word for this chapter.
\end{grammarlens}

\subsection*{Your turn}
\begin{itemize}
\raggedright
  \item \emph{Say it:} \emph{gôlā}
  \item \emph{Say it:} \emph{gôlā}, once more
  \item \emph{Say it:} say \emph{gôlā}
  \item \emph{Recall:} say the Bengali for hair, then the Bengali for a tooth, then say \emph{gôlā} again
\end{itemize}

\begin{tcolorbox}[breakable,colback=teal!4,colframe=teal!35!black,title={Before you move on}]
What does \bn{গলা} mean? (\textquotedblleft{}The throat\textquotedblright{}.) Say it once more.
\end{tcolorbox}
\section[\texorpdfstring{\bn{পেট} (peṭ)}{পেট (peṭ)}]{\bn{পেট} (peṭ) — the stomach}

\label{lesson:BN-C42-pet}



\begin{quote}
Before the new one: say the Bengali for a tooth, then the Bengali for the throat.
\end{quote}

\subsection*{You'll want to know: \bn{পেট}}
\textbf{\bn{পেট}} — \emph{peṭ} — \textquotedblleft{}the stomach\textquotedblright{}.

\begin{cousinweb}
\textbf{\bn{পেট}} (\emph{peṭ}) is a tadbhava, a word inherited and reshaped by sound change. It comes from \emph{*peṭṭa}, \textquotedblleft{}belly\textquotedblright{}, a form CDIAL reconstructs from its descendants (CDIAL 8376). CDIAL notes that it may be one word with a similar \emph{*pēṭṭa}, \textquotedblleft{}basket\textquotedblright{}.
\end{cousinweb}

\begin{grammarlens}[title={What you've built}]
One more word for this chapter.
\end{grammarlens}

\subsection*{Your turn}
\begin{itemize}
\raggedright
  \item \emph{Say it:} \emph{peṭ}
  \item \emph{Say it:} \emph{peṭ}, once more
  \item \emph{Say it:} say \emph{peṭ}
  \item \emph{Recall:} say the Bengali for a tooth, then the Bengali for the throat, then say \emph{peṭ} again
\end{itemize}

\begin{tcolorbox}[breakable,colback=teal!4,colframe=teal!35!black,title={Before you move on}]
What does \bn{পেট} mean? (\textquotedblleft{}The stomach\textquotedblright{}.) Say it once more.
\end{tcolorbox}
